CMAT provides additional mathematics support that students may use during the academic period. Each visit is recorded in the administrative system, and throughout this paper a \emph{visit} refers to one recorded instance of attendance at CMAT. The records identify whether and how often a student attended, but they do not measure the duration or content of the support received, the amount of independent study that followed a visit, or the nature of the student's interaction with support staff. The number of visits is therefore interpreted as the frequency of recorded attendance, not as a measure of tutoring time.

CMAT attendance during MU also takes place within an institutional incentive structure. During part of the study period, three recorded visits were relevant to a first-semester participation requirement. This is reported because it may influence the distribution of attendance, but the present paper does not evaluate that requirement and does not use the three-visit value as a discontinuity or causal threshold.

The administrative record may contain CMAT visits from more than one academic period for the same student. Using cumulative attendance would therefore assign visits made before or after MU to the grade from a particular MU attempt. To avoid this mismatch, the analysis counts only visits recorded in the same academic period as the MU attempt being analysed.

\section{Data and methods}

\subsection{Analytical population}

The primary analytical population contains one observation per student, corresponding to the first eligible attempt at MU. Administrative records that do not represent a course attempt, such as equivalencies and revalidations, are excluded from the attempt definition, whereas adverse academic outcomes are retained in the primary outcome construction. The analysis is restricted to academic periods for which CMAT attendance records are available, because students with no recorded visits can only be treated as non-attenders when the attendance system covers the corresponding period. There are 7,777 eligible first-MU attempts across all observed periods, of which 6,627 occur during periods with CMAT attendance coverage and form the primary analytical cohort.

The primary cohort spans 190 instructor--period classes. Class size has a median of 30 students (IQR 24--48), a mean of 34.9, and a range of 10--88. Of the 190 classes, 176 contain at least one student who attended CMAT and at least one student with no recorded visits. For the specific attendance-group comparisons, 152 classes contain both zero- and one-visit students, 114 contain both zero- and two-visit students, 84 contain both zero- and three-visit students, and 102 contain both zero- and four-or-more-visit students. Thus, the fixed-effect comparisons are supported by substantial within-class variation in attendance, although fewer classes contribute to comparisons involving the smaller attendance groups.

\subsection{Attendance}

For each student, attendance is measured as the number of CMAT visits recorded during the same academic period as the MU attempt. The main analysis uses the five attendance groups shown in \Cref{eq:mu-visit-groups}:
\begin{equation}
V\in\{0,\;1,\;2,\;3,\;4+\}.
\label{eq:mu-visit-groups}
\end{equation}
This grouping distinguishes students with no recorded attendance from students who attended once, twice, three times, or at least four times, without assuming a linear relationship between visit frequency and performance. The four-or-more category is used because individual visit counts become increasingly sparse at higher levels of attendance.

\subsection{Outcome and class standardisation}

The primary outcome is the final MU grade standardised within class, as defined in \Cref{eq:classroom-relative-performance}. For student $i$ in class $c$,
\begin{equation}
Z_{ic}=\frac{Y_{ic}-\bar{Y}_{c}}{s_c},
\label{eq:classroom-relative-performance}
\end{equation}
where a class is defined as students taught by the same instructor in the same course and academic period. Because the primary analysis includes only MU, this is effectively an instructor-by-period grouping. The resulting score expresses each student's performance relative to other students whose work was evaluated in the same local grading context.

Standardising grades within class is a measurement choice rather than a claim that all differences between instructors or periods are error. Raw course grades may reflect differences in assessment, grading practice, cohort composition, or other class-level conditions. Within-class standardisation reduces the extent to which differences in grade distributions across classes can drive a pooled association between CMAT attendance and grades, but it does not address student-level self-selection into mathematics support. The resulting $Z$ score should therefore be interpreted as classroom-relative academic performance, not as a general measure of mathematical proficiency.

Non-numeric adverse outcomes are retained in the primary specification using the study's prespecified imputation procedure before within-class standardisation. Because the treatment of these outcomes can affect students' relative positions, we also report a sensitivity analysis using a common imputed value for adverse non-numeric outcomes and a complete-case specification restricted to numeric final grades.

\subsection{Statistical analysis}

We first report group-specific means and 95\% confidence intervals for $Z$ across the five attendance groups. Because the Brown--Forsythe test indicates unequal variances across groups, equality of means is assessed using Welch's one-way ANOVA. Pairwise post hoc comparisons use the Games--Howell procedure, which accommodates unequal group sizes and variances. These analyses describe differences in academic performance across attendance groups; they do not establish that attendance at mathematics support caused those differences.

